% ===== PRÉAMBULE CANONIQUE — à coller VERBATIM en tête de CHAQUE .tex =====
\documentclass[11pt,a4paper]{article}
\usepackage[utf8]{inputenc}\usepackage[T1]{fontenc}\usepackage{lmodern}
\usepackage[french]{babel}
\usepackage{amsmath,amssymb,mathtools}\usepackage{icomma}
\usepackage[version=4]{mhchem}          % \ce{...} : équations & formules
\usepackage{siunitx}\sisetup{locale=FR} % \qty{}{}, \si{}, \ang{}
\usepackage{chemfig}                    % \chemfig{...} : structures développées
\usepackage{xcolor}\usepackage{colortbl}
\usepackage{tikz}\usepackage{pgfplots}\pgfplotsset{compat=1.18}
\usepackage[most]{tcolorbox}\usepackage{enumitem}\usepackage{array,booktabs}
\usepackage[a4paper,margin=2.2cm]{geometry}
\usetikzlibrary{arrows.meta,calc,angles,quotes,positioning,patterns,backgrounds,shapes.geometric}
\definecolor{primary}{RGB}{222,49,99}\definecolor{primarydark}{RGB}{150,22,55}
\definecolor{defcol}{RGB}{0,114,178}\definecolor{methcol}{RGB}{52,92,148}
\definecolor{excol}{RGB}{0,135,90}\definecolor{attncol}{RGB}{200,40,55}
\tcbset{boxbase/.style={enhanced,breakable,boxrule=0.9pt,arc=2.5pt,
  left=3mm,right=3mm,top=2.3mm,bottom=2.3mm,fonttitle=\bfseries,coltitle=white}}
% --- boîtes TUTORIEL (noms EXACTS, ne pas renommer) ---
\newtcolorbox{cle}[1][]{boxbase,colback=defcol!6,colframe=defcol,
  title={\textsc{À retenir}\ifx\relax#1\relax\relax\else\textmd{ : \itshape #1}\fi}}
\newtcolorbox{methode}[1][]{boxbase,colback=methcol!7,colframe=methcol,sharp corners,
  title={\textsc{Méthode}\ifx\relax#1\relax\relax\else\textmd{ --- \itshape #1}\fi}}
\newtcolorbox{exemplebox}[1][]{boxbase,colback=excol!5,colframe=excol,
  title={\textsc{Exemple}\ifx\relax#1\relax\relax\else\textmd{ : \itshape #1}\fi}}
\newtcolorbox{piege}[1][]{boxbase,colback=attncol!6,colframe=attncol,
  title={\textsc{Piège}\ifx\relax#1\relax\relax\else\textmd{ : \itshape #1}\fi}}
\newtcolorbox{remarque}[1][]{boxbase,colback=black!4,colframe=black!45,
  title={\textsc{Remarque}\ifx\relax#1\relax\relax\else\textmd{ : \itshape #1}\fi}}
% --- boîtes EXERCICE / CORRIGÉ (noms EXACTS, auto-comptées) ---
\newtcolorbox[auto counter]{exo}[1][]{boxbase,colback=primary!4,colframe=primary,
  title={\textsc{Exercice}~\thetcbcounter\ifx\relax#1\relax\relax\else\textmd{ --- \itshape #1}\fi}}
\newtcolorbox[auto counter]{corrige}[1][]{boxbase,colback=excol!4,colframe=excol,
  title={\textsc{Corrigé}~\thetcbcounter\ifx\relax#1\relax\relax\else\textmd{ --- \itshape #1}\fi}}
\newcommand{\R}{\mathbb{R}}\newcommand{\N}{\mathbb{N}}\newcommand{\Z}{\mathbb{Z}}\newcommand{\ds}{\displaystyle}
% =========================================================================
\begin{document}
\sisetup{per-mode=symbol}
\begin{corrige}[établir la relation par le tableau d'avancement]
\begin{enumerate}[label=\alph*)]
  \item \ce{Ag3PO4} : $[\ce{Ag+}]=3s$ et $[\ce{PO4^{3-}}]=s$, donc
    \[ K_{\mathrm{ps}} = [\ce{Ag+}]^{3}[\ce{PO4^{3-}}] = (3s)^{3}(s) = 27s^{3}\times s = 27\,s^{4}. \]
  \item \ce{Ca3(PO4)2} : $[\ce{Ca^{2+}}]=3s$ et $[\ce{PO4^{3-}}]=2s$, donc
    \[ K_{\mathrm{ps}} = [\ce{Ca^{2+}}]^{3}[\ce{PO4^{3-}}]^{2} = (3s)^{3}(2s)^{2}
       = 27s^{3}\times 4s^{2} = 108\,s^{5}. \]
\end{enumerate}
\end{corrige}

\begin{corrige}[calculer $s$ à partir de $K_{\mathrm{ps}}$ (types variés)]
On isole $s$ dans l'expression du type correspondant.
\begin{enumerate}[label=\alph*)]
  \item $s = \sqrt{K_{\mathrm{ps}}} = \sqrt{\num{3.3e-9}} = \qty{5.7e-5}{\mole\per\litre}$.
  \item $s = \sqrt[3]{\dfrac{K_{\mathrm{ps}}}{4}} = \sqrt[3]{\dfrac{\num{6.6e-6}}{4}}
           = \sqrt[3]{\num{1.65e-6}} = \qty{1.18e-2}{\mole\per\litre}$.
  \item $s = \sqrt[4]{\dfrac{K_{\mathrm{ps}}}{27}} = \sqrt[4]{\dfrac{\num{2.0e-39}}{27}}
           = \sqrt[4]{\num{7.4e-41}} \approx \qty{9.3e-11}{\mole\per\litre}$.
\end{enumerate}
\end{corrige}

\begin{corrige}[$K_{\mathrm{ps}}$ à partir d'une variable imposée]
La concentration $[\ce{Mn^{2+}}]$ \emph{est} la solubilité : $s=[\ce{Mn^{2+}}]=\dfrac{x}{2}$. Par la
stœchiométrie ($1{:}2$), $[\ce{OH-}]=2s=2\times\dfrac{x}{2}=x$. Donc
\[ K_{\mathrm{ps}} = [\ce{Mn^{2+}}][\ce{OH-}]^{2} = \left(\dfrac{x}{2}\right)(x)^{2} = \dfrac{x^{3}}{2}. \]
(Cohérent avec le type $1{:}2$ : $K_{\mathrm{ps}}=4s^{3}=4(x/2)^{3}=x^{3}/2$.)
\end{corrige}

\begin{corrige}[classer par solubilité (même type)]
Les trois sels étant de \emph{même type} ($1{:}1$, $s=\sqrt{K_{\mathrm{ps}}}$), la solubilité croît avec
$K_{\mathrm{ps}}$. On peut classer directement par $K_{\mathrm{ps}}$ croissant :
\[ \underbrace{\ce{BaSO4}}_{\num{1.1e-10}} \ <\ \underbrace{\ce{CaC2O4}}_{\num{2.3e-9}}
   \ <\ \underbrace{\ce{CaCO3}}_{\num{3.3e-9}}. \]
Vérification par $s=\sqrt{K_{\mathrm{ps}}}$ : $s(\ce{BaSO4})=\qty{1.0e-5}{}$,
$s(\ce{CaC2O4})=\qty{4.8e-5}{}$, $s(\ce{CaCO3})=\qty{5.7e-5}{\mole\per\litre}$ : même ordre.
\end{corrige}

\begin{corrige}[de la solubilité massique au $K_{\mathrm{ps}}$]
On convertit d'abord $s_m$ en $s$, puis on applique $K_{\mathrm{ps}}=4s^{3}$ (type $1{:}2$).
\[ s = \dfrac{s_m}{M} = \dfrac{\qty{4.84}{\gram\per\litre}}{\qty{367.0}{\gram\per\mole}}
     = \qty{1.32e-2}{\mole\per\litre}, \]
\[ K_{\mathrm{ps}} = 4s^{3} = 4\,(\num{1.32e-2})^{3} = 4\times\num{2.30e-6} = \num{9.2e-6}. \]
\end{corrige}
\end{document}
